\chapter{Symptoms of serious illness in babies and children}
\index{Symptoms of serious illness in babies and children}

\medskip
\noindent\textit{Educational reference; always seek professional advice for individual care.}

\section*{Definition and Overview}
Babies and children can become seriously unwell quickly, and recognising the early warning signs is one of the most valuable skills a parent or carer can develop. Serious illness means a condition with the potential to progress rapidly to serious harm, such as sepsis, meningitis, or severe dehydration. Children are not simply small adults: their immune systems are still developing, their bodies lose fluid quickly, and their early symptoms are often subtle or non-specific. A baby who is "just a bit quiet" may in fact be seriously unwell, which is why experienced clinicians place as much weight on how a child looks and behaves as on the height of a fever or any single symptom.

Fever is extremely common in childhood and is usually caused by a harmless viral infection, but in a baby under three months of age it is always taken seriously. Most serious childhood illness is infectious, but injuries, congenital conditions, and complications of chronic disease such as asthma or diabetes also matter. This chapter provides an overview of the red-flag symptoms that should prompt urgent assessment, the main conditions they point to, and what typically happens when a child is taken to a doctor or hospital.

A key principle throughout is trust in parental instincts: parents are usually right when they feel something is seriously wrong, even when no specific symptom stands out. local health professionals are trained to listen carefully when a parent says their child is "not right". The sections that follow describe how to assess a child at home, the patterns that suggest particular serious illnesses, when to call an ambulance, and how dangerous conditions are investigated and treated.

\section*{Epidemiology}
The overwhelming majority of acute illness in children is caused by viral infection and resolves without treatment, but a small minority of unwell children have a serious bacterial infection. The most important of these are sepsis, meningitis, pneumonia, and urinary tract infection. Among young infants with fever, serious bacterial infection is uncommon, affecting in the order of one in every 1000 to 2000 febrile babies, but it is sufficiently dangerous that doctors investigate it carefully. Children under five years of age have the highest rates of acute illness, and boys are slightly more often affected than girls for many infections.

Patterns have changed dramatically over recent decades. Immunisation has almost eliminated the common causes of bacterial meningitis and serious pneumonia, although cases still occur, particularly in the unvaccinated, and early antibiotics and better hospital care have improved outcomes. In many countries, children in remote areas and Aboriginal and Torres Strait Islander communities experience higher rates of some severe infections, and this remains a focus of public health work. These patterns help families understand that serious illness is rare but real, and that vaccination and early medical review are the most powerful protections.

\section*{Aetiology and Pathophysiology}
The commonest serious illnesses of childhood are infections, and they harm the body through several mechanisms. Bacterial and viral pathogens can cause systemic inflammation that progresses to sepsis, in which the blood pressure falls and organs begin to fail; they can spread to the brain and its lining to cause meningitis; they can inflame the airways and lungs to cause croup, bronchiolitis, and pneumonia; and they can cause vomiting and diarrhoea that drain the body of fluid faster than a small child can replace it. Non-infectious causes also matter, including congenital heart disease, metabolic disorders, injuries, and surgical emergencies such as intussusception, in which part of the bowel telescopes into itself.

Infants are vulnerable for several reasons: they cannot tell anyone how they feel, their immune systems are still maturing, and their bodies have small fluid and temperature reserves. A baby's airway is narrow and easily obstructed by inflammation, their skull allows the fontanelle to bulge if pressure rises inside the brain, and a fever in the first weeks of life can be the only sign of a serious infection.

\begin{itemize}
    \item \textbf{Immature immunity:} young infants produce a weaker and less targeted immune response, so bacteria can multiply and spread more readily
    \item \textbf{Rapid fluid loss:} children have proportionally larger fluid losses with vomiting and diarrhoea, and dehydrate faster than adults
    \item \textbf{Limited communication:} infants and toddlers cannot describe pain, headache, or feeling ill, so behaviour change is often the only clue
    \item \textbf{Anatomical vulnerability:} narrow airways, a soft larynx, and the open fontanelle shape how symptoms present
    \item \textbf{Incomplete vaccination:} a child who is unvaccinated or behind schedule is at higher risk of vaccine-preventable serious illness
\end{itemize}

\section*{Clinical Features}
The features below are the red flags that should prompt urgent medical assessment. A child can have one or several of them, and not all serious illnesses produce all of them.

\begin{itemize}
    \item \textbf{Fever:} any fever in a baby under three months, or a fever that persists beyond three days, needs review; very high or rising temperatures are taken seriously in the first year
    \item \textbf{Breathing difficulty:} fast or noisy breathing, grunting, pulling in of the chest or ribs, flaring of the nostrils, or a child who cannot talk in full sentences
    \item \textbf{Lethargy:} hard to wake, unusually floppy, staring blankly, or "not themselves" without obvious reason
    \item \textbf{Poor feeding:} a baby taking far less than usual, refusing drinks, or having very few wet nappies
    \item \textbf{Non-blanching rash:} a rash that does not fade when pressed with a glass, particularly with fever
    \item \textbf{Seizure or fit:} a convulsion, stiffness with shaking, or a period of blank unresponsiveness
    \item \textbf{Bulging fontanelle or stiff neck:} a tense soft spot, stiff neck, or high-pitched cry can indicate meningitis
    \item \textbf{Bluish lips or tongue:} blue or very pale, mottled skin suggests poor oxygen delivery
    \item \textbf{Inconsolable crying:} a baby who cries constantly and cannot be settled, or who arches their back when held
    \item \textbf{Dehydration signs:} sunken eyes, dry mouth, cold hands and feet, and few wet nappies
\end{itemize}

\section*{Differential Diagnosis}
The red flags above overlap with a range of conditions, and distinguishing them drives the medical assessment.

\begin{itemize}
    \item \textbf{Common viral infection:} the most frequent cause of fever in children; the child is usually alert, drinking, and improving within a few days
    \item \textbf{Sepsis:} suspected when a feverish child looks toxic, is lethargic, or has poor perfusion; requires urgent blood tests and antibiotics
    \item \textbf{Meningitis:} suggested by stiff neck, bulging fontanelle, high-pitched cry, and a non-blanching rash; needs immediate assessment
    \item \textbf{Pneumonia and bronchiolitis:} respiratory viruses and bacteria cause fast breathing and chest signs; more dangerous in infants
    \item \textbf{Urinary tract infection:} a common cause of unexplained fever in infants and toddlers, confirmed by urine testing
    \item \textbf{Gastroenteritis and dehydration:} vomiting and diarrhoea with reduced urine output and lethargy
    \item \textbf{Febrile seizure:} a brief convulsion with fever, usually harmless but needing review to exclude meningitis
    \item \textbf{Intussusception:} intermittent severe abdominal pain, drawing up of the legs, and vomiting
    \item \textbf{Congenital heart disease:} poor feeding, rapid breathing, and cyanosis may be the first clues in early infancy
\end{itemize}

\section*{When to Seek Medical Care}
Any child who appears seriously unwell should be assessed promptly. As a rule, trust your instincts: if a baby or child looks or behaves very differently from usual, seek medical advice even if you cannot name the problem. See a GP or child health service for a fever in a baby under three months, a fever lasting more than three days, poor feeding, reduced wet nappies, or a child who is listless but able to be roused. Seek advice earlier if a child has a chronic condition, is under six months, or has been in contact with serious illness.

\textbf{Contact your local emergency services for:}
\begin{itemize}
    \item Difficulty breathing, such as grunting, chest indrawing, or flaring nostrils
    \item Blue or very pale, mottled lips, tongue, or skin
    \item A seizure, especially one lasting more than five minutes
    \item Unconsciousness or a child who cannot be roused
    \item A non-blanching rash with fever
    \item Severe dehydration with sunken eyes, a dry mouth, and no wet nappies for many hours
    \item A bulging fontanelle, stiff neck, or high-pitched cry
\end{itemize}

\section*{Diagnosis and Investigations}
When a seriously unwell child reaches medical care, the assessment is designed to be rapid and safe, and usually begins with checking breathing, circulation, and responsiveness before any history is taken.

\begin{itemize}
    \item \textbf{History:} the clinician asks about onset, fever, feeding, wet nappies, breathing, behaviour, vaccination, and contact with illness
    \item \textbf{Physical examination:} temperature, breathing rate and effort, oxygen saturation, skin and rash, fontanelle, neck, ears, throat, and abdomen
    \item \textbf{Urine testing:} a sample is checked for urinary tract infection in febrile infants and toddlers
    \item \textbf{Blood tests:} full blood count, inflammatory markers, and blood cultures are used when sepsis is suspected
    \item \textbf{Chest X-ray:} used when pneumonia or heart problems are possible
    \item \textbf{Lumbar puncture:} spinal fluid is sampled if meningitis is suspected, before or shortly after antibiotics
    \item \textbf{Abdominal imaging:} ultrasound or X-ray can confirm intussusception or other surgical emergencies
\end{itemize}

\section*{Management}
Treatment depends on the underlying condition and the child's age, but the principles below apply in most situations.

\begin{itemize}
    \item \textbf{Reassurance and observation:} most unwell children have viral infections and need rest, fluids, and time
    \item \textbf{Fever control:} paracetamol and ibuprofen are used to keep a child comfortable; aspirin is never used in children
    \item \textbf{Fluids:} encourage small, frequent drinks; oral rehydration solution is useful for gastroenteritis; children who cannot drink may need intravenous fluids
    \item \textbf{Antibiotics:} given promptly, often intravenously, when bacterial infection such as sepsis or meningitis is suspected
    \item \textbf{Oxygen and breathing support:} used for bronchiolitis, pneumonia, or respiratory illness with low oxygen levels
    \item \textbf{Hospital admission:} infants under three months with fever, and any child with suspected sepsis or meningitis, is usually admitted
    \item \textbf{Treating the cause:} specific treatment, such as surgery for intussusception or insulin for diabetic emergencies
    \item \textbf{Follow-up:} arrange review within a day or two for children managed at home, with clear instructions on what should prompt an earlier return
\end{itemize}

\section*{Complications}
Most children with a serious illness recover fully, but several complications can occur, especially when treatment is delayed.

\begin{itemize}
    \item \textbf{Febrile seizures:} convulsions with fever are frightening but usually leave no lasting harm
    \item \textbf{Dehydration:} can progress to poor circulation and shock if fluid losses are not replaced
    \item \textbf{Sepsis and septic shock:} organ failure can follow overwhelming infection and requires intensive care
    \item \textbf{Meningitis:} can cause hearing loss, brain injury, learning difficulties, or limb loss in severe cases
    \item \textbf{Pneumonia:} fluid around the lungs or lung abscess in severe cases
    \item \textbf{Delayed diagnosis:} the greatest avoidable harm comes from missing a serious illness while it is still treatable
\end{itemize}

\section*{Prognosis and Outcomes}
The outlook for seriously ill children is very good when care is sought early. Bacterial sepsis and meningitis, once commonly fatal, now carry a low risk of death in many countries because of rapid antibiotics, intensive care, and vaccination. Febrile seizures do not cause brain damage, and pneumonia usually resolves fully with prompt treatment. The key to a good outcome is timing: the earlier a serious infection is recognised and treated, the fewer the complications.

Parents play the central role in that timing. Knowing the red flags, acting on instinct, and seeking review when a child is "not right" save lives, and it is far better to have a child assessed and reassured than to wait at home with a progressing condition. For children with ongoing needs after a serious illness, follow-up with the GP, paediatrician, and allied health services supports full recovery and normal development.

\section*{Prevention and Self-Care}
\begin{itemize}
    \item Keep immunisation up to date; it is the most powerful protection against serious bacterial illness
    \item Practise good hand hygiene, especially around babies and before preparing feeds
    \item Keep babies and children away from cigarette smoke, which increases the risk and severity of respiratory infection
    \item Follow safe-sleep practices for babies, including sleeping on the back in a safe cot
    \item Use age-appropriate car seats and supervise children around water and roads
    \item Learn the red flags in this chapter and keep the healthdirect or local emergency numbers handy
    \item Trust your instincts and seek medical review early if a child is unwell and "not themselves"
\end{itemize}

\section*{Sources and Further Reading}
The following organisations provide reliable information about serious illness in children for parents and clinicians.

\begin{itemize}
    \item NHS. \textit{Fever in babies and children: when to see a doctor.} https://www.nhs.uk/
    \item Raising Children Network. \textit{When your child is sick: what to watch for.} https://raisingchildren.net.au/
    \item Mayo Clinic. \textit{Symptoms of serious illness in infants and children.} https://www.mayoclinic.org/
    \item healthdirect Australia. \textit{When to take a baby or child to the emergency department.} https://www.healthdirect.gov.au/
    \item World Health Organization (WHO). \textit{Caring for a sick child in the community.} https://www.who.int/
    \item Australian Department of Health and Aged Care. \textit{Childhood immunisation schedule and information.} https://www.health.gov.au/
\end{itemize}
\clearpage
